\documentclass[11pt]{article}
\usepackage[margin=1in]{geometry}
\usepackage{amsmath,amssymb,amsthm}
\usepackage{graphicx}
\usepackage[hidelinks]{hyperref}

\newtheorem{theorem}{Theorem}
\newtheorem{lemma}{Lemma}
\newtheorem{remark}{Remark}
\newcommand{\Z}{\mathbb Z}
\newcommand{\R}{\mathbb R}
\newcommand{\PhiF}{\Phi}
\newcommand{\sk}{\mathsf{s}_{k,\Phi,w}}
\newcommand{\skp}{\mathsf{s}_{k',\Phi,w}}
\newcommand{\lk}{l_{k,\Phi,w}}
\newcommand{\lkp}{l_{k',\Phi,w}}

\title{A Counterexample to Compactness Without the Zero-\(\Delta_2\) Condition}
\author{Agent lane 14, run \texttt{fa\_banach\_001}}
\date{2026-07-03}

\begin{document}
\maketitle

\section*{Source and Target}

The source paper is Pierre-A. Vuillermot,
\emph{On the compactness of embeddings for a class of weighted
Orlicz-Sobolev sequence spaces}, arXiv:2503.19454.  The target is the remark
after Theorem 4 on PDF page 17, source lines 1188--1195.  In the notation of
the source paper, the author asks whether the compact embedding
\[
  \mathsf{s}_{k',\Phi,w}\hookrightarrow \mathsf{s}_{k,\Phi,w},\qquad k'>k,
\]
still holds when \(\mathsf{s}_{k,\Phi,w}\) is strictly included in
\(l_{k,\Phi,w}\), and also whether compact embeddings
\[
  l_{k',\Phi,w}\hookrightarrow l_{k,\Phi,w},\qquad k'>k,
\]
exist in that situation.

\begin{center}
\includegraphics[width=0.88\linewidth]{figures/open_problem_crop.png}
\end{center}

\section*{Proof Intuition}

The source proof of compactness uses the zero-\(\Delta_2\) condition to make
tails uniformly small after passing from the \(k'\)-weight to the \(k\)-weight.
If \(\Phi\) is very flat at the origin, this uniform tail estimate can fail
exactly on one-coordinate vectors.  Choosing \(\Phi(t)=e^{-1/t}\) near zero
makes the inverse behave like \(1/\log(1/u)\).  For the far-out coordinate
\(n\), the \(k'\)-unit vector has size
\[
  \frac{1}{k'\log(1+\Phi(n))},
\]
and its \(k\)-norm is therefore the constant \(k/k'\), independent of \(n\).
These vectors cannot have a compact image.

\section*{Counterexample}

\begin{theorem}
For every pair \(0<k<k'\), there are an Orlicz function \(\Phi\) and weights
\(w_m\equiv 1\) such that:
\begin{enumerate}
\item \(\mathsf{s}_{k,\Phi,w}\subsetneq l_{k,\Phi,w}\);
\item the embedding \(\mathsf{s}_{k',\Phi,w}\to \mathsf{s}_{k,\Phi,w}\) is
continuous but not compact;
\item the embedding \(l_{k',\Phi,w}\to l_{k,\Phi,w}\) is continuous but not
compact.
\end{enumerate}
Consequently the compactness questions in the source remark have a negative
answer in general.
\end{theorem}

\begin{proof}
Fix \(0<k<k'\).  Let \(0<a<1/2\), for instance \(a=1/4\), and define
\[
\Phi(0)=0,\qquad \Phi(t)=e^{-1/t}\quad (0<t\leq a),
\]
with a linear extension for \(t\geq a\):
\[
\Phi(t)=\Phi(a)+\Phi'(a)(t-a),\qquad t\geq a.
\]
Since \(e^{-1/t}\) has second derivative
\[
  e^{-1/t}\frac{1-2t}{t^4}
\]
on \((0,a]\), this function is convex on \([0,a]\).  The linear extension
uses the tangent slope at \(a\), so the full function is continuous, convex,
increasing, positive on \((0,\infty)\), and unbounded.  Thus it is an Orlicz
function in the sense of the source paper.  It fails the zero-\(\Delta_2\)
condition because
\[
  \frac{\Phi(2t)}{\Phi(t)}
  =\exp\!\left(\frac{1}{2t}\right)\to\infty
  \qquad (t\downarrow 0).
\]
Set \(w_m=1\) for all \(m\in\Z\).

We first show that \(\mathsf{s}_{k,\Phi,w}\) is strictly smaller than
\(l_{k,\Phi,w}\).  For all sufficiently large positive integers \(m\), set
\[
  p_m=\frac{1}{(k+2)\log(1+\Phi(m))},
\]
and set \(p_m=0\) for the remaining indices.  Since the extension of \(\Phi\)
is linear at infinity, \(1+\Phi(m)\) is comparable to \(m\).  Also \(p_m<a\)
for large \(m\).  Hence
\[
  \Phi(p_m)=(1+\Phi(m))^{-(k+2)}
\]
for large \(m\), and therefore
\[
  \sum_m (1+\Phi(|m|))^k\Phi(|p_m|)
  <\infty.
\]
Thus \(p=(p_m)\in l_{k,\Phi,w}\).  On the other hand, take
\(\rho=(k+1)/(k+2)\).  Then, for large \(m\),
\[
  \Phi(p_m/\rho)=(1+\Phi(m))^{-(k+1)},
\]
so
\[
  \sum_m (1+\Phi(|m|))^k\Phi(|p_m|/\rho)
  \geq \sum_{m\gg 1}(1+\Phi(m))^{-1}=\infty,
\]
again because \(\Phi(m)\) is linear at infinity.  Thus \(p\notin
\mathsf{s}_{k,\Phi,w}\).

Now define, for large positive integers \(n\),
\[
  x^{(n)}=\frac{1}{k'\log(1+\Phi(n))}\,e_n,
\]
where \(e_n\) is the \(n\)-th coordinate vector.  These vectors have finite
support, hence belong to every small space \(\mathsf{s}_{j,\Phi,w}\).  Put
\(U_n=1+\Phi(n)\).  Since the nonzero coordinate is in the
\(e^{-1/t}\)-range of \(\Phi\), we have
\[
  \Phi\!\left(\frac{1}{k'\log U_n}\right)=U_n^{-k'}.
\]
The Luxemburg norm of a one-coordinate vector \(b e_n\) in \(l_{j,\Phi,w}\)
is
\[
  \|b e_n\|_{j,\Phi,w}
  =\frac{|b|}{\Phi^{-1}((1+\Phi(n))^{-j})}.
\]
Consequently
\[
  \|x^{(n)}\|_{k',\Phi,w}=1,
  \qquad
  \|x^{(n)}\|_{k,\Phi,w}=\frac{k}{k'}.
\]
Thus \((x^{(n)})\) is contained in the unit ball of \(\mathsf{s}_{k',\Phi,w}\)
and of \(l_{k',\Phi,w}\), but its image in the \(k\)-space has constant norm
\(k/k'\).

For each coordinate \(n\), the coordinate projection \(P_n\) is contractive
for the Luxemburg norm, since deleting coordinates can only decrease the
defining modular.  Hence, if \(n\ne m\),
\[
  \|x^{(n)}-x^{(m)}\|_{k,\Phi,w}
  \geq \|P_n(x^{(n)}-x^{(m)})\|_{k,\Phi,w}
  =\|x^{(n)}\|_{k,\Phi,w}
  =\frac{k}{k'}.
\]
The image sequence has no Cauchy subsequence in \(l_{k,\Phi,w}\), and a
fortiori no Cauchy subsequence in the closed subspace
\(\mathsf{s}_{k,\Phi,w}\).  Therefore neither embedding is compact.  The
continuity of both embeddings for \(k'>k\) is the elementary continuous
embedding already recorded in the source paper before Theorem 3.
\end{proof}

\section*{Verification Notes}

No numerical computation is used.  The key checks are:
\begin{itemize}
\item the function \(\Phi\) is a valid Orlicz function and fails the
zero-\(\Delta_2\) condition;
\item the sequence \(p_m=[(k+2)\log(1+\Phi(m))]^{-1}\) proves strict
inclusion \(\mathsf{s}_{k,\Phi,w}\subsetneq l_{k,\Phi,w}\);
\item the one-coordinate norm formula gives the constant image norm \(k/k'\);
\item contractive coordinate projections prevent any Cauchy subsequence.
\end{itemize}

\section*{Novelty Check}

The run indexes were searched for arXiv:2503.19454, the source title, compact
embeddings, and the exact notation \(\mathsf{s}_{k,\Phi,w}\),
\(l_{k,\Phi,w}\); no duplicate packet or prior attempt was found.  A bounded
web search on 2026-07-03 for exact title/id phrases and the displayed
compactness notation found the source arXiv page but no later answer to this
exact question.  This packet should therefore be reviewed as a new candidate
counterexample, not as a literature-status identification.

\section*{References}

\begin{itemize}
\item P.-A. Vuillermot, \emph{On the compactness of embeddings for a class of
weighted Orlicz-Sobolev sequence spaces}, arXiv:2503.19454, 2025.
\end{itemize}

\end{document}
